% !TEX root = ../manuscript.tex

\section*{Note sur l'utilisation de l'intelligence artificielle}
\addcontentsline{toc}{chapter}{Note sur l'utilisation de l'IA}

Dans le cadre de ce mémoire, j'ai eu recours à des outils d'intelligence artificielle générative --- principalement \textbf{Claude} (Anthropic) --- comme aide à la recherche et à la rédaction. Voici comment ces outils ont été utilisés concrètement~:

\begin{itemize}
  \item \textbf{Recherche bibliographique} : exploration et synthèse de la littérature sur les thématiques abordées (machine learning en santé, méthodes de régularisation, gestion du déséquilibre de classes, interprétabilité).
  \item \textbf{Structuration du document} : aide à l'organisation des chapitres et à la formulation des transitions.
  \item \textbf{Relecture et reformulation} : suggestions de reformulation pour améliorer la clarté des explications techniques.
  \item \textbf{Débogage et code} : aide au débogage des notebooks Python (scikit-learn, XGBoost, SHAP) et à la compilation \LaTeX.
\end{itemize}

L'ensemble du contenu scientifique --- choix des modèles, interprétation des résultats, analyse critique, conclusions --- a été compris, vérifié et validé par l'auteur. L'IA a servi d'outil d'assistance~; elle n'a pas remplacé le raisonnement ni le travail analytique qui constituent le c\oe ur de ce mémoire. Toutes les erreurs éventuelles restent de ma seule responsabilité.

\bigskip
